\chapter*{B3\quad《流体力学与流体机械》试卷 3（网络题库截图）}
\addcontentsline{toc}{chapter}{B3\quad《流体力学与流体机械》试卷 3（网络题库截图）}
\markboth{B3\quad《流体力学与流体机械》试卷 3（网络题库截图）}{}

\begin{tishi}
\textbf{来源：}本卷是一张\textbf{网络文库（百度文库）网页截图}，题名《流体力学与流体机械》试卷 3。
原文档为在线文档的 6 屏截图，\textbf{题干、选项与文库页面的侧栏推荐文档、阅读数、按钮等杂讯交错}，
本册已逐题从杂讯中剥离出来，\textbf{题号与选项一律按剥题结果重新排列}，
题干中的待填处统一记作（\quad）。

\textbf{题型与分值（原文档原文）：}单项选择题 30 题 $\times$ 2 分 $=$ 60 分；
多项选择题 10 题 $\times$ 2 分 $=$ 20 分；判断题 20 题 $\times$ 1 分 $=$ 20 分，合计 100 分。
\textbf{原文档没有计算题}，故本卷只出上述三部分。

\textbf{剥题订正说明（逐条）：}
① 第 1 题原文档作“牛顿流体是（\quad）的一种物质”，但其四个选项都在解释\emph{流体}的定义、
标准答案也是流体的定义，故应为“\textbf{流体是（\quad）的一种物质}”，本册按后者订正；
② 第 9 题原文档识别一度读作“可压缩流体恒定总流的连续性方程”，
但四个选项与标准答案都是\emph{不可压缩}流体总流连续性方程的形式，本册订为“\textbf{不可压缩流体}”；
③ 第 12 题选项 D 原文档印为“减少阀门的启闭时间”（截图上另有手写“延长”二字，属读者批注，非原题文字），
本册照录原印刷文字；
④ 第 24 题原文档题干“比离心式泵与风机的流量（\quad）”文字完整，未作改动；
⑤ 原文档“粘滞”统一作“\textbf{黏滞}”、“素流”订为“\textbf{紊流}”、“蝇壳”订为“\textbf{蜗壳}”，
微分符号统一为 $\dd$，题设数值与选项数值一律未作改动。

\textbf{超纲提示：}本卷涉及\textbf{泵与风机（流体机械）}的题共 22 题
（单选第 21--30 题，多选第 4、5、6、10 题，判断第 1--6、8、10 题）。
流体机械对应赵琴教材第 11--14 章，\textbf{不在 818 初试范围}（初试只考前 8 章），
本册一律标最低档 \xstarfull{1}，并在依据里注明“超纲：流体机械（泵与风机）……复试范围”。
\end{tishi}

\section*{一、单项选择题（30 题 $\times$ 2 分，共 60 分）}

\begin{ti}
\sloppy
\emergencystretch=2em

\item 流体是（\quad）的一种物质。
\begin{choices}
\item 不断膨胀直到充满容器的
\item 实际上是不可压缩的
\item 不能承受剪切力的
\item 在任一剪切力的作用下不能保持静止的
\end{choices}
\xstarfull{4}\yiju{E1 slide33“第一章绪论＝重点掌握基础概念”；E1 slide16“选择/判断题考概念辨析”；E5 本题库单选第 1 题}\nandu{易}

\item 连续性方程表示（\quad）。
\begin{choices}
\item 能量守恒
\item 动量守恒
\item 质量守恒
\item 动量矩守恒
\end{choices}
\xstarfull{5}\yiju{E1 slide33“第四章流体动力学基础＝考试计算题重点”；E2“第四章是主要考点，需掌握计算套路”；E5 本题库单选第 2、9 题}\nandu{易}

\item 层流的沿程阻力系数（\quad）。
\begin{choices}
\item 仅与雷诺数有关
\item 仅与相对粗糙度有关
\item 与雷诺数及相对粗糙度有关
\item 是常数
\end{choices}
\xstarfull{5}\yiju{E1 slide33“第五章孔口/管嘴/管路水力计算＝考试计算题重点”；E2“第五章同为主要计算考点”；E5 本题库单选第 3、15 题}\nandu{易}

\item 动量方程中流速和作用力（\quad）。
\begin{choices}
\item 流速有方向，作用力没有方向
\item 流速没有方向，作用力有方向
\item 都没有方向
\item 都有方向
\end{choices}
\xstarfull{5}\yiju{E1 slide33“第四章＝考试计算题重点”；E4 2015 年试题计算第 4 题（水平弯管动量方程求 $F_x$、$F_y$）；E4 2026 年回忆版计算第 3 题（教材动量方程习题 4.12）}\nandu{中}

\item 雷诺数是表征（\quad）。
\begin{choices}
\item 惯性力与黏滞力之比
\item 惯性力与重力之比
\item 水压力与黏滞力之比
\item 黏滞力与重力之比
\end{choices}
\xstarfull{4}\yiju{E4 2015 年试题填空（$\Rey$ 判流态）、2022 年单选第 9 题（两圆管流态判别）、2026 年回忆版简答（层流与雷诺数关系）；E1 slide33“第八章＝选择填空题考点”}\nandu{易}

\item 圆管水流的下临界雷诺数为（\quad）。
\begin{choices}
\item 200
\item 2000
\item 4000
\item 10000
\end{choices}
\xstarfull{4}\yiju{E4 2015 年试题填空（$\Rey$ 判流态）、2022 年单选第 9 题（两圆管流态判别）；E1 slide33“第八章＝选择填空题考点”}\nandu{易}

\item 按长管计算的串联管道，无支流分出，则各管段的（\quad）。
\begin{choices}
\item 流量相等
\item 切应力相等
\item 水力坡度相等
\item 沿程水头损失相等
\end{choices}
\xstarfull{5}\yiju{E1 slide33“第五章＝考试计算题重点”；E2“第五章同为主要计算考点”；E5 本题库单选第 7、13 题（长管的串联与并联计算）}\nandu{易}

\item 某流体的运动黏度 $\nu=3\times10^{-6}\ \mathrm{m^2/s}$、密度 $\rho=800\ \mathrm{kg/m^3}$，其动力黏度 $\mu$ 为（\quad）。
\begin{choices}
\item $3.75\times10^{-9}\ \mathrm{Pa\cdot s}$
\item $2.4\times10^{-3}\ \mathrm{Pa\cdot s}$
\item $2.4\times10^{5}\ \mathrm{Pa\cdot s}$
\item $2.4\times10^{9}\ \mathrm{Pa\cdot s}$
\end{choices}
\xstarfull{4}\yiju{E4 2015 年试题填空第 1 题（重度与动力黏度换算）；E1 slide33“第一章绪论＝重点掌握基础概念”}\nandu{易}

\item 不可压缩流体恒定总流的连续性方程是（\quad）。
\begin{choices}
\item $u_1\dd A_1=u_2\dd A_2$
\item $u_1A_1=u_2A_2$
\item $v_1A_1=v_2A_2$
\item $v_1\dd A_1=v_2\dd A_2$
\end{choices}
\xstarfull{5}\yiju{E1 slide33“第四章＝考试计算题重点”；E2“第四章是主要考点，需掌握计算套路”；E5 本题库单选第 2 题}\nandu{易}

\item 底宽 $b=2\ \mathrm{m}$、水深 $h=1\ \mathrm{m}$ 的矩形断面渠道，水力半径 $R$ 等于（\quad）。
\begin{choices}
\item $0.5\ \mathrm{m}$
\item $1.0\ \mathrm{m}$
\item $1.5\ \mathrm{m}$
\item $2.0\ \mathrm{m}$
\end{choices}
\xstarfull{2}\yiju{E6 教材第 5 章非圆管水力半径 $R=A/\chi$（管路沿程损失计算的基础量）；E5 本题库单选第 10 题；E4 历年真题未直接出现，故压至两星}\nandu{易}

\item 流线与迹线一定重合的流动是（\quad）。
\begin{choices}
\item 无旋流动
\item 均匀流动
\item 平面流动
\item 定常流动
\end{choices}
\xstarfull{4}\yiju{E4 2015 年试题计算第 1 题（流线方程与流动方向判断）、第 6 题（流函数求速度与流线）；E1 slide33“第三章流体运动学＝公式多、重点记忆”}\nandu{中}

\item 为减轻管道中的水击压强，不能采取的措施是（\quad）。
\begin{choices}
\item 选用直径较大的管道
\item 选用管壁较薄的管道
\item 减小管道长度
\item 减少阀门的启闭时间
\end{choices}
\xstarfull{4}\yiju{E4 2026 年回忆版简答考“水击现象”；E1 slide33“第五章＝考试计算题重点”；E5 本题库单选第 12 题}\nandu{中}

\item 按长管计的并联管道，各支管的（\quad）。
\begin{choices}
\item 流量相等
\item 切应力相等
\item 沿程水头损失相等
\item 水力坡度相等
\end{choices}
\xstarfull{5}\yiju{E1 slide33“第五章＝考试计算题重点”；E2“第五章同为主要计算考点”；E5 本题库单选第 7、13 题}\nandu{易}

\item 仅受重力时，静止液体的测压管水头线为（\quad）。
\begin{choices}
\item 水平线
\item 斜线
\item 平行线
\item 铅垂线
\end{choices}
\xstarfull{4}\yiju{E4 2022 年单选第 3 题（测压管水头线坡度小于零时的沿程变化）；E1 slide33“第二章流体静力学＝计算题基础入门、重点掌握方法”}\nandu{易}

\item 圆管紊流过渡区的沿程水头损失系数 $\lambda$ 与（\quad）有关。
\begin{choices}
\item 雷诺数
\item 管壁相对粗糙度
\item 雷诺数和管壁相对粗糙度
\item 雷诺数和管道长度
\end{choices}
\xstarfull{5}\yiju{E1 slide33“第五章＝考试计算题重点”；E4 2015 年试题计算第 2 题（突然缩小的水头损失）；E5 本题库单选第 3 题}\nandu{中}

\item 某水库水深为 $10\ \mathrm{m}$ 处的绝对压强为（\quad）。
\begin{choices}
\item $10\ \mathrm{kPa}$
\item $98\ \mathrm{kPa}$
\item $147\ \mathrm{kPa}$
\item $196\ \mathrm{kPa}$
\end{choices}
\xstarfull{4}\yiju{E1 slide33“第二章＝计算题基础入门、重点掌握方法”（静压强基本计算）；E4 2022 年单选第 7 题（虹吸管最高处压强与大气压的关系）}\nandu{易}

\item 下列关于流体压强的说法中，正确的是（\quad）。
\begin{choices}
\item 绝对压强可正可负
\item 相对压强可正可负
\item 真空压强可正可负
\item 绝对压强恒为正，而真空压强恒为负
\end{choices}
\xstarfull{4}\yiju{E4 2022 年单选第 7 题（绝对压强、相对压强与大气压的关系）；E1 slide33“第二章＝重点掌握方法”}\nandu{易}

\item 给定某一瞬间，如果流场中每一空间点上流体相对应的物理参数均相等，这种流场称为（\quad）。
\begin{choices}
\item 定常流场
\item 非定常流场
\item 均匀流场
\item 非均匀流场
\end{choices}
\xstarfull{3}\yiju{E1 slide16“简答题考……均匀流特点等”；E1 slide33“第三章＝公式多、重点记忆”；E5 本题库单选第 18 题}\nandu{易}

\item 当圆管中流体作层流流动时，动能修正系数 $\alpha$ 等于（\quad）。
\begin{choices}
\item 1
\item 2
\item 3
\item 2000
\end{choices}
\xstarfull{5}\yiju{E1 slide33“第四章＝考试计算题重点”（伯努利方程中的动能修正系数）；E4 2015 年试题填空（层流断面平均速度等于 0.5 倍最大速度，$\alpha=2$ 的直接来源）}\nandu{中}

\item 流体在管内作层流流动时，其沿程损失 $h_f$ 值与断面平均流速 $v$ 的（\quad）次方成正比。
\begin{choices}
\item 1
\item 1.75
\item $1.75\sim2$
\item 2
\end{choices}
\xstarfull{5}\yiju{E1 slide33“第五章＝考试计算题重点”；E1 slide16“简答题考水头损失分类及定义”；E5 本题库单选第 3、15 题}\nandu{中}

\item 平衡盘在多级离心泵中的作用是平衡（\quad）。
\begin{choices}
\item 径向力
\item 轴向力
\item 液流冲击力
\item 惯性力
\end{choices}
\xstarfull{1}\yiju{超纲：流体机械（泵与风机），赵琴第 11--14 章，复试范围}\nandu{易}

\item 下列能自动调整轴向力平衡的装置是（\quad）。
\begin{choices}
\item 平衡盘
\item 平衡鼓
\item 平衡管
\item 平衡叶片
\end{choices}
\xstarfull{1}\yiju{超纲：流体机械（泵与风机），赵琴第 11--14 章，复试范围}\nandu{易}

\item 叶轮与泵壳间用的密封是（\quad）密封。
\begin{choices}
\item 软填料
\item 硬填料
\item 口环
\item 机械
\end{choices}
\xstarfull{1}\yiju{超纲：流体机械（泵与风机），赵琴第 11--14 章，复试范围}\nandu{易}

\item 一般情况下，轴流式泵与风机比离心式泵与风机的流量（\quad）。
\begin{choices}
\item 小
\item 大
\item 相等
\item 不能比较大小
\end{choices}
\xstarfull{1}\yiju{超纲：流体机械（泵与风机），赵琴第 11--14 章，复试范围}\nandu{易}

\item 泵与风机的实际工作点应落在（\quad）点附近，工作才最经济。
\begin{choices}
\item 最大压头
\item 最大功率
\item 最高效率
\item 最大流量
\end{choices}
\xstarfull{1}\yiju{超纲：流体机械（泵与风机），赵琴第 11--14 章，复试范围（与 E5 题库计算题第 4 题“水泵工况点与节流调节”同类）}\nandu{易}

\item 叶轮的作用是使流体获得（\quad）。
\begin{choices}
\item 动能
\item 压能
\item 能量
\item 速度
\end{choices}
\xstarfull{1}\yiju{超纲：流体机械（泵与风机），赵琴第 11--14 章，复试范围}\nandu{易}

\item 按工作原理，叶片式泵与风机一般为轴流式、混流式和（\quad）。
\begin{choices}
\item 滑片式
\item 螺杆式
\item 往复式
\item 离心式
\end{choices}
\xstarfull{1}\yiju{超纲：流体机械（泵与风机），赵琴第 11--14 章，复试范围}\nandu{易}

\item 轴流式风机的工作原理是利用旋转叶轮、叶片对流体作用的（\quad）来输送流体，并提高其压力。
\begin{choices}
\item 压力
\item 升力
\item 离心力
\item 重力
\end{choices}
\xstarfull{1}\yiju{超纲：流体机械（泵与风机），赵琴第 11--14 章，复试范围}\nandu{易}

\item 风机蜗壳的作用是（\quad）。
\begin{choices}
\item 导向流体
\item 使流体加速
\item 使流体的能量增加
\item 收集流体，并使流体的部分动能转变为压能
\end{choices}
\xstarfull{1}\yiju{超纲：流体机械（泵与风机），赵琴第 11--14 章，复试范围}\nandu{易}

\item 蜗壳形状是按液流（\quad）这一自由轨迹而设计的。
\begin{choices}
\item 螺旋线
\item 双曲线
\item 抛物线
\item 曲线
\end{choices}
\xstarfull{1}\yiju{超纲：流体机械（泵与风机），赵琴第 11--14 章，复试范围}\nandu{易}

\end{ti}

\section*{二、多项选择题（10 题 $\times$ 2 分，共 20 分）}

\begin{ti}
\sloppy
\emergencystretch=2em

\item 描述流体运动的方法有（\quad）。
\begin{choices}
\item 欧拉法
\item 拉格朗日法
\item 分析法
\item 以上都是
\end{choices}
\xstarfull{3}\yiju{E1 slide33“第三章＝公式多、重点记忆”；E1 slide16“选择/判断题考概念辨析”；E5 本题库多选第 1 题}\nandu{易}

\item 决定流体流态的雷诺数 $\Rey$ 与（\quad）有关。
\begin{choices}
\item 流体的比重
\item 流体的平均流速
\item 运动黏滞系数
\item 流体的静压力
\item 管子的直径
\end{choices}
\xstarfull{4}\yiju{E4 2015 年试题填空（$\Rey$ 判流态）、2022 年单选第 9 题（两圆管流态判别）、2026 年回忆版简答（层流与雷诺数关系）；E1 slide33“第八章＝选择填空题考点”}\nandu{中}

\item 描写流体运动状态的运动参数有（\quad）。
\begin{choices}
\item 平均速度
\item 重力加速度
\item 速度
\item 加速度
\item 角速度
\end{choices}
\xstarfull{3}\yiju{E1 slide33“第三章＝公式多、重点记忆”；E1 slide33“第四章＝考试计算题重点”（运动参数是列方程的前提）；E5 本题库多选第 3 题}\nandu{中}

\item 叶片式流体机械中的能量损失可分为（\quad）。
\begin{choices}
\item 水力损失
\item 容积损失
\item 机械损失
\item 以上都是
\end{choices}
\xstarfull{1}\yiju{超纲：流体机械（泵与风机），赵琴第 11--14 章，复试范围}\nandu{易}

\item 水泵固定部分主要有（\quad）。
\begin{choices}
\item 吸水段
\item 排水段
\item 轴承
\item 中段
\end{choices}
\xstarfull{1}\yiju{超纲：流体机械（泵与风机），赵琴第 11--14 章，复试范围}\nandu{易}

\item 200D43$\times$3 型泵型号的意义是（\quad）。
\begin{choices}
\item 吸入口直径 200 mm
\item 单吸多级分段式
\item 单级额定扬程为 43 m
\item 3 级
\end{choices}
\xstarfull{1}\yiju{超纲：流体机械（泵与风机），赵琴第 11--14 章，复试范围}\nandu{易}

\item 绝对压强 $p_{abs}$ 与相对压强 $p$、真空度 $p_v$、当地大气压 $p_a$ 之间的关系是（\quad）。
\begin{choices}
\item $p_{abs}=p_v+p_a$
\item $p=p_{abs}-p_a$
\item $p_v=p_a-p_{abs}$
\item $p=p_{abs}+p_a$
\end{choices}
\xstarfull{4}\yiju{E4 2022 年单选第 7 题（绝对压强、相对压强与大气压的关系）；E1 slide33“第二章＝计算题基础入门、重点掌握方法”}\nandu{中}

\item 根据管道水头损失计算方法的不同，管道可以分为（\quad）。
\begin{choices}
\item 长管
\item 短管
\item 圆管
\item 粗糙管
\item 光滑管
\end{choices}
\xstarfull{5}\yiju{E1 slide33“第五章＝考试计算题重点”；E2“第五章同为主要计算考点”；E5 本题库单选第 7、13 题（长管串联与并联计算）}\nandu{易}

\item 流体的物理特性有密度、比容和（\quad）。
\begin{choices}
\item 体积
\item 重量
\item 黏性
\item 压缩性
\item 重度
\end{choices}
\xstarfull{4}\yiju{E4 2015 年试题填空第 1 题（重度与动力黏度换算）；E1 slide33“第一章＝重点掌握基础概念”}\nandu{中}

\item 通风机噪声的控制方法有（\quad）。
\begin{choices}
\item 降低声源的噪声
\item 传输路径的噪声控制
\item 接收器的防护
\item 以上都是
\end{choices}
\xstarfull{1}\yiju{超纲：流体机械（泵与风机），赵琴第 11--14 章，复试范围}\nandu{易}

\end{ti}

\section*{三、判断题（20 题 $\times$ 1 分，共 20 分）}

\begin{ti}
\sloppy
\emergencystretch=2em

\item （\quad）如果泵几何相似，则比转数相等下的工况为相似工况。
\xstarfull{1}\yiju{超纲：流体机械（泵与风机）及泵的相似换算，赵琴第 11--14 章，复试范围}\nandu{中}

\item （\quad）按流体吸入叶轮的方式可将离心泵分为轴流泵和混流泵。
\xstarfull{1}\yiju{超纲：流体机械（泵与风机），赵琴第 11--14 章，复试范围}\nandu{易}

\item （\quad）离心泵启动时排液管上的闸阀必须关闭，以减小启动功率。
\xstarfull{1}\yiju{超纲：流体机械（泵与风机），赵琴第 11--14 章，复试范围}\nandu{易}

\item （\quad）按级数可将离心泵分为单吸式泵和双吸式泵。
\xstarfull{1}\yiju{超纲：流体机械（泵与风机），赵琴第 11--14 章，复试范围}\nandu{易}

\item （\quad）容积式泵主要包括往复泵和离心泵。
\xstarfull{1}\yiju{超纲：流体机械（泵与风机），赵琴第 11--14 章，复试范围}\nandu{易}

\item （\quad）泵的运行工况点是流量--扬程曲线和流量--效率曲线的交点。
\xstarfull{1}\yiju{超纲：流体机械（泵与风机），赵琴第 11--14 章，复试范围}\nandu{易}

\item （\quad）相对压强可以大于、等于或小于零。
\xstarfull{4}\yiju{E4 2022 年单选第 7 题（虹吸管最高处压强与大气压的关系，须分清相对压强与真空压强）；E1 slide33“第二章＝重点掌握方法”}\nandu{易}

\item （\quad）泵与风机在管路系统中工作时，必须满足能量的供求平衡。
\xstarfull{1}\yiju{超纲：流体机械（泵与风机），赵琴第 11--14 章，复试范围}\nandu{易}

\item （\quad）有压长管道是认为水头全部消耗在沿程水头损失上。
\xstarfull{5}\yiju{E1 slide33“第五章＝考试计算题重点”；E2“第五章同为主要计算考点”；E5 本题库单选第 7、13 题（长管的串联与并联计算）}\nandu{易}

\item （\quad）水泵的扬程就是指水泵的提水高度。
\xstarfull{1}\yiju{超纲：流体机械（泵与风机），赵琴第 11--14 章，复试范围}\nandu{易}

\item （\quad）总流连续性方程是从能量守恒方程推导出来的。
\xstarfull{5}\yiju{E1 slide33“第四章＝考试计算题重点”；E2“第四章是主要考点，需掌握计算套路”；E5 本题库单选第 2、9 题}\nandu{易}

\item （\quad）紊流边界层是指边界层内流体以紊流运动为主。
\xstarfull{3}\yiju{E1 slide33“第八章黏性流体动力学基础＝选择填空题考点”；E1 slide16“选择/判断题考概念辨析”；E5 本题库判断第 12 题}\nandu{中}

\item （\quad）水力光滑管是指层流底层厚度大于管壁粗糙度、粗糙峰对流动无影响的流动。
\xstarfull{4}\yiju{E1 slide33“第五章＝考试计算题重点”（沿程阻力系数 $\lambda$ 的分区判断）；E1 slide16“选择/判断题考概念辨析”；E5 本题库单选第 15 题}\nandu{中}

\item （\quad）对于不可压缩液体，动力黏度相等的液体，其运动黏度也相等。
\xstarfull{3}\yiju{E4 2015 年试题填空第 1 题（运动黏度与动力黏度的换算 $\nu=\mu/\rho$）；E1 slide33“第一章＝重点掌握基础概念”}\nandu{中}

\item （\quad）体积模量 $K$ 值越大，液体越容易压缩。
\xstarfull{3}\yiju{E1 slide33“第一章＝重点掌握基础概念”（压缩性与体积模量）；E5 题库（B2）第一章第 8 题}\nandu{易}

\item （\quad）无论是静止的还是运动的液体，都不存在切力。
\xstarfull{4}\yiju{E1 slide33“第一章＝重点掌握基础概念”（黏性只在运动时显现，静止流体不承受切应力）；E1 slide16“选择/判断题考概念辨析”}\nandu{易}

\item （\quad）液体在静止状态下和运动状态下都能具有黏滞性。
\xstarfull{4}\yiju{E1 slide33“第一章＝重点掌握基础概念”；E5 题库（B1）问答题第 1 题“流动性与黏滞性”}\nandu{易}

\item （\quad）尼古拉兹试验是研究管道沿程水头损失随雷诺数和相对粗糙度的变化关系的试验。
\xstarfull{4}\yiju{E1 slide33“第五章＝考试计算题重点”（沿程阻力的分区与尼古拉兹试验）；E5 本题库单选第 3、15 题}\nandu{中}

\item （\quad）在常温、常压下水的运动黏度比空气的运动黏度大。
\xstarfull{3}\yiju{E1 slide33“第一章＝重点掌握基础概念”（运动黏度 $\nu=\mu/\rho$ 及其随温度的变化）；E4 2015 年试题填空第 1 题}\nandu{易}

\item （\quad）两种不同种类的液体，只要流速梯度相等，它们的切应力也相等。
\xstarfull{4}\yiju{E1 slide33“第一章＝重点掌握基础概念”（牛顿内摩擦定律 $\tau=\mu\dfrac{\dd u}{\dd y}$）；E5 题库（B1）问答题第 1 题“流动性与黏滞性”}\nandu{易}

\end{ti}
